\فصل{مقدمه}
\برچسب{chap:intro}

فرآیند \واژه{تست.ن.ا.} از اولین پیدایش تا کنون با رویکرد‌ها و کاربردهای
متنوعی اجرا شده‌است. در این پایان‌نامه بر رویکرد
بهره‌گیری از توصیف نرم‌افزار به زبان مشتری در \واژه{تست} تمرکز شده و
معماری‌ جدیدی را برای \واژه{تست} با \واژه{تکنیک} \واژه{بی.دی.دی.}، که بر اساس این
رویکرد ارائه شده، پیشنهاد می‌کند.

گرچه ایده‌ی مطرح شده در معماری پیشنهادی منحصر به  \واژه{تکنیک}
\واژه{بی.دی.دی.} نیست و با سایر  \واژه{تکنیک}‌ها و \واژه{متدولوژی}‌ها نیز قابل انطباق است،
حوزه‌ی بررسی‌این پایان‌نامه به \واژه{تکنیک} \واژه{بی.دی.دی.} محدود شده
است.

در بخش \رجوع{sec:intro:history} تاریخچه‌ای از رویکردهای موجود نسبت به
\واژه{تست.ن.ا.} ارائه می‌شود. در بخش ‌های \رجوع{sec:intro:tdd} و
\رجوع{sec:intro:bdd} به ترتیب ادبیات مورد نیاز و جزییات \واژه{تکنیک}
\واژه{بی.دی.دی.} شرح داده می‌شوند.


\قسمت{تاریخچه}\برچسب{sec:back:history}

روند تکامل رویکرد رایج در \واژه{تست.ن.ا.}، از سال ۱۹۵۷ تا سال ۲۰۰۰ به
پنج دوره‌ی \نام{رفع‌اشکال-محور}{Debugging-oriented}،
\نام{اثبات-محور}{Demonstration-oriented}،
\نام{تخریب-محور}{Destruction-oriented}،
\نام{ارزیابی-محور}{Evaluation-oriented} و
\نام{پیشگیری-محور}{Prevention-oriented} تقسیم‌بندی
‌می‌شود\مرجع{laycock1993theory}. برخی فعالیت‌های بارز این دوره‌ها به شرح
زیر است: \شروع{شمارش}

\فقره{ تا سال ۱۹۵۷، برنامه نویس‌ها پس از نوشتن برنامه، آن را اجرا نموده
تا از کارکرد صحیح آن اطمینان یابند و اگر \واژه{باگ}ی در اجرا پیدا
می‌شد، آن را پیدا کرده و رفع می‌نمودند. این فرآیند ادامه پیدا می‌کرد تا
زمانی که احساس کنند \واژه{باگ}ی باقی نمانده است.  }

\فقره{ در سال ۱۹۵۷، \نام{بیکر}{Charles L. Baker} برای اولین بار
\واژه{تست.ن.ا.} را، به‌عنوان روشی برای «اثبات عملکرد صحیح» برنامه، از
\واژه{دیباگ} تمیز داد\مرجع{baker1957review}.  }

\فقره{ \نام{دایسترا}{Edsger Dijkstra} در سال ۱۹۶۹ متذکر شد کاربرد تست،
«یافتن \واژه{باگ}ها» است، نه اثبات عملکرد صحیح
برنامه\مرجع{testinghistory}.  در سال ۱۹۸۳، راهنمایی برای
\واژه{اعتبار.سنجی}، \واژه{وارسی} و \واژه{تست.ن.ا.} منتشر شد
\مرجع{adrion1982validation} که «تشخیص \واژه{باگ}‌های تحلیل و طراحی» را،
علاوه بر تشخیص \واژه{باگ}‌های پیاده‌سازی، با \واژه{تست.ن.ا.} میسر می‌سازد
\مرجع{luo2001software}.  }

\فقره{ در سال ۱۹۸۷، معیار \واژه{اتکا.پذیری} به‌عنوان عاملی کلیدی در
«اندازه‌گیری \واژه{کیفیت.ن.ا.}» معرفی شد \مرجع{musa1987software}.  }

\فقره{ در سال ۱۹۹۰، \واژه{تست.ن.ا.} به عنوان یکی از موثرترین عوامل
«پیشگیری از \واژه{باگ}»‌ معرفی شد \مرجع{beizer20672software}.  }

\پایان{شمارش}

پس از رویکردهای مذکور در اواخر دهه ۱۹۹۰، با ظهور \واژه{متدولوژی}‌های
جدید مانند \واژه{ایکس.پی.}، که در دسته‌بندی \واژه{متدولوژی}‌های
\واژه{ای.اس.دی.} قرار می‌گیرد، رویکردهای جدیدی در \واژه{تست.ن.ا.} ایجاد
شد.

در \واژه{متدولوژی}‌های سنتی با \واژه{مدل.آبشاری}، \واژه{تست} یکی از
فاز‌های \واژه{اس.دی.پی.} بود که فقط یک بار انجام می‌شد و پیش‌نیاز آن،
اتمام فازهای \واژه{تحلیل}، \واژه{طراحی} و \واژه{پیاده.سازی} بود. در
\واژه{متدولوژی}‌های آبشاری، هر کدام از فازهای \واژه{تحلیل}،
\واژه{طراحی}، \واژه{پیاده.سازی} و \واژه{تست} به‌صورت مجزا اجرا می‌شد
\مرجع{agileVsWaterfall} و \واژه{تست} نقش موثری در فازهای تحلیل، طراحی
و پیاده‌سازی نداشت.

اما در \واژه{متدولوژی}‌های چابک، مراحل تحلیل، طراحی، پیاده سازی و
\واژه{تست} به‌صورت \واژه{تکراری} اجرا می‌شود. در هر \واژه{تکرار}،
نیازمندی کامل نرم‌افزار مشخص نیست و در \واژه{تکرار} بعد، ممکن است تغییر
کند. نقش \واژه{تست} در «پاسخگویی به تغییرات»، که یکی از چهار ارزش
\واژه{بیانیه.اجایل} است \مرجع{manifesto}، بسیار کلیدی است. در
\واژه{تست} چابک، رویکرد «\نام{دستیاری در کیفیت}{Quality Assistance}»
بر رویکرد «\نام{اطمینان از کیفیت}{Quality Assurance}» غلبه می‌کند\مرجع{ghahrai}.

\نام{بک}{Kent Beck} در سال ۲۰۰۲، با معرفی \واژه{تکنیک}
\واژه{تی.دی.دی.}، \واژه{تست} را به‌عنوان «محرک توسعه‌ی نرم‌افزار» معرفی
می‌کند و نوشتن \واژه{تست} را پیش‌نیاز \واژه{پیاده.سازی} می‌داند. وی
«بهبود طراحی نرم‌افزار» را از نتایج بکارگیری این \واژه{تکنیک} می‌داند
\مرجع{beck2003test}. همچنین با بهره‌گیری از این نگرش در \واژه{متدولوژی}
\واژه{ای.ام.دی.دی.}، از \واژه{تست.ن.ا.} به‌عنوان «\واژه{توصیف}
نرم‌افزار» بهره گرفته می‌شود \مرجع{agileModeling:specByExample}.


%%% Local Variables:
%%% mode: latex
%%% TeX-master: "../main"
%%% End:

\زیرقسمت{\واژه{تی.دی.دی.}}\برچسب{sec:back:tdd}

\واژه{تی.دی.دی.} \واژه{تکنیک}ی است که برنامه‌نویس را به استفاده از دو
\واژه{تکنیک} \واژه{تی.اف.دی.} و \واژه{ریفکتورینگ} با دو هدف اصلی زیر
ملزم می‌کند \مرجع{agiledata:tdd}:

\شروع{شمارش}

\فقره { با بهره‌گیری از \واژه{تی.اف.دی.}، برنامه‌نویس موظف است قبل از
  پیاده‌سازی هر واحد از کد، به نیازمندی آن واحد و طراحی آن واحد فکر
  کرده و در حقیقت با نوشتن \واژه{تست}، نیازمندی آن واحد را نیز توصیف
  می‌کند. }

\فقره { با توجه به اینکه برنامه‌نویس موظف است پس از افزودن هر واحد،
  \واژه{ریفکتورینگ} را به‌عنوان مرحله‌ای اجباری از فرآیند توسعه‌ی
  نرم‌افزار لحاظ کند، تمیزی کد کیفیت طراحی نرم‌افزار همواره در بهترین
  حالت حفظ می‌شود. }

\پایان{شمارش}

همانطور که در نمودار حالت شکل \رجوع{fig:tdd-states} نمایش داده شده‌است،
در حالت \واژه{سبز}، فقط امکان \واژه{ریفکتورینگ} هست و برای تغییر متن
برنامه، باید ابتدا با نوشتن یک \واژه{تست} برای آن، به حالت
\واژه{قرمز} رفته و سپس امکان تغییر متن برنامه را داریم.


\شروع{شکل}[tbp]
\تنظیم‌ازوسط
\درج‌تصویر[پهنا=0.7\پهنای‌سطر]{tdd-states}
\شرح
[نمودار حالت وضعیت‌های \واژه{تست.سویت} و نحوه‌ی گذار بین آنها]
{نمودار حالت وضعیت‌های \واژه{تست.سویت} و نحوه‌ی گذار بین آنها با
  استفاده افزودن \واژه{تست}، \واژه{ریفکتورینگ} و تغییر متن برنامه \مرجع{agiledata:tdd}}
\برچسب{fig:tdd-states}
\پایان{شکل}


%%% Local Variables:
%%% mode: latex
%%% TeX-master: "../main"
%%% End:

\قسمت{\واژه{تکنیک} \واژه{بی.دی.دی.}}\برچسب{sec:back:bdd}

\واژه{بی.دی.دی.} یکی از تکنیک‌های توسعه‌ی چابک نرم‌افزار است که در سال
۲۰۰۶ توسط \نام{نورث}{Dan North} برای پاسخ‌گویی به ابهام‌ها و مشکل‌هایی که
در \واژه{تی.دی.دی.} بوجود آمده‌بود، معرفی
شد. \مرجع{north2006introducing}.  تاکید \واژه{بی.دی.دی.} بر بهینه‌سازی
ارتباط میان نقش‌های برنامه‌نویس، آزمون‌گر و \واژه{دامین.اکسپرت}‌ می‌باشد و
برای تیمی که تمام اعضای آن برنامه‌نویس باشند، \واژه{بی.دی.دی.} سودی
نسبت به \واژه{تی.دی.دی.} ندارد \مرجع{north2012}.

این تکنیک، که اخیرا هم در عمل و هم در پژوهش بسیار متداول شده است، حاصل
ترکیب و بهبود \واژه{پرکتیس}‌های معرفی شده در تکنیک‌های \واژه{تی.دی.دی.}
و \واژه{ای.تی.دی.دی.} و علاوه‌بر آن در نظر گرفتن اصول زیر است
\مرجع{alience:bdd}:

\شروع{بارهها}

\باره{ لازم است هدف هر \واژه{یوزر.استوری} و سودی که برای مشتری دارد،
  مشخص باشد.}

\باره{ با نگرش «\نام{از بیرون به درون}{From the outside in}»، فقط
  رفتارهایی از \واژه{سیستم} پیاده‌سازی می‌شوند که بیشترین سود را به
  مشتری می‌رسانند. با داشتن این نگرش، اتلاف هزینه کمینه می‌شود. }

\باره{ رفتارهای مورد انتظار از سامانه، بین \واژه{دامین.اکسپرت}،
  توسعه‌دهنده‌ی \واژه{سیستم} و آزمون‌گر به یک زبان مشترک توصیف می‌شوند. در نتیجه
  ارتباط ضروری میان این نقش‌ها بهبود می‌یابد. }

\باره { این اصول، در تمام سطوح انتزاعی توصیف برنامه، تا پایین‌ترین
  سطح که پیاده‌سازی یک واحد است، رعایت می‌شوند. }

\پایان{بارهها}

\زیرقسمت{مشخصات \واژه{بی.دی.دی.}}  در حال حاضر، \واژه{تکنیک}
\واژه{بی.دی.دی.}، هنوز در حال توسعه است و تعریف روشنی از
\واژه{بی.دی.دی.} که بر آن اجماع باشد، وجود ندارد. با توجه به اینکه
فرآیند \واژه{بی.دی.دی.}، از ابتدا به‌صورت انتزاعی معرفی شده و جزییاتی
برای آن ارائه نشده، مشخصات \واژه{بی.دی.دی.} مبهم و پراکنده
هستند. همچنین چارچوب‌ها و ابزارهایی که برای \واژه{بی.دی.دی.} ارائه
شده‌اند، بیشتر به بخش «پیاده‌سازی آزمون» از فرآیند \واژه{بی.دی.دی.}
تمرکز دارند؛ در صورتی که \واژه{بی.دی.دی.} بر حوزه‌ی گسترده‌تری از
\واژه{اس.دی.پی.} تاثیر دارد \مرجع{solis2011study}.

در مقاله‌ی \مرجع{solis2011study}، ۶ مورد از مشخصات اصلی
\واژه{بی.دی.دی.} که بر کل \واژه{اس.دی.پی.}، نه فقط قسمت «پیاده‌سازی
آزمون»، تاثیر گذار اند، ارائه شده است. در اینجا به صورت خلاصه به آن‌ها
اشاره می‌کنیم:

\شروع{شمارش}

\فقره{ \متن‌سیاه{زبان \واژه{یوبی}}:

  مفهوم «زبان \واژه{یوبی}»، هسته‌ی \واژه{بی.دی.دی.} را تشکیل
  می‌دهد. ساختار \واژه{یوبی} برآمده از \واژه{بیزنس.دامین} می‌باشد و
  ابهام را از مکالمه‌ی بین مشتری و تیم توسعه کاهش می‌دهد. همچنین یک
  واژه‌نامه در ابتدای پروژه ایجاد شده، اکثر واژه‌های آن در مرحله‌ی تحلیل
  افزوده می‌شوند و در مراحل بعد امکان گسترش دارد.

  هر \واژه{بیزنس.دامین}، زبان \واژه{یوبی} خاص خود را لازم دارد.
  \واژه{بی.دی.دی.} یک \واژه{قالب} ساده برای مرحله‌ی تحلیل ارائه نموده است که
  مستقل از \واژه{بیزنس.دامین} است و از آن در زبان \واژه{یوبی} بهره‌گیری
  می‌شود. این \واژه{قالب} در بخش \رجوع{gwt} شرح داده شده است.  }

\فقره{ \متن‌سیاه{فرآیند \واژه{تفکیک} \واژه{تکراری}}:

  توقع مشتری از یک پروژه‌ی نرم‌افزاری، ایجاد \واژه{بیزنس.ولیو}
  می‌باشد. معمولاً تشخیص و روشن‌سازی \واژه{بیزنس.ولیو}، دشوار است. به
  همین دلیل، \واژه{بیزنس.ولیو} به اجزای ملموس‌تر تفکیک می‌شود. در شکل
  \رجوع{fig:bdd-conceptual} رابطه‌ی اجزای استفاده شده در این
  \واژه{تفکیک} نسبت به هم نمایش داده شده است.

  \شروع{شکل}[tbp]
  \تنظیم‌ازوسط
  \درج‌تصویر[پهنا=1\پهنای‌سطر]{bdd-conceptual}
  \تر\موقتم{
    رابطه‌ی اجزای تشکیل‌دهنده‌ی \واژه{بیزنس.ولیو}
  }
  \شرح[\موقتم]{\موقتم~\مرجع{solis2011study}}
  \برچسب{fig:bdd-conceptual}
  \پایان{شکل}
 
  مرحله‌ی تحلیل در \واژه{بی.دی.دی.}، با شناسایی \واژه{رفتار}های مورد
  انتظار از سامانه آغاز می‌شود. در مراحل ابتدایی، شناسایی رفتارهای مورد
  انتظار آسان‌تر از شناسایی ارزش‌های تجاری سامانه است. هر رفتار تجاری،
  تعدادی \واژه{نتیجه.تجاری} محسوس را محقق می‌سازد.

  هر \واژه{نتیجه.تجاری} با یک \واژه{فیچر.ست} اعمال می‌شود. یک
  \واژه{فیچر.ست}‌ با گفتگو بین تیم توسعه و مشتری است برای تحقق
  \واژه{نتیجه.تجاری} در سامانه تدوین گشته و اولویت‌بندی و تبیین صریح
  ارتباظ آن با \واژه{نتیجه.تجاری}، ضروری است.

  در نهایت \واژه{اسکوپ} هر \واژه{فیچر.ست}، با استفاده از چند
  \واژه{یوزر.استوری} معین می‌شود. هر واژه{یوزر.استوری}، یک تعامل بین
  کاربر و سامانه را توصیف می‌کند. هر \واژه{یوزر.استوری}، به سه
  سوال زیر پاسخ می‌دهد:

  \شروع{بارهها}

  \باره{نقش کاربر در \واژه{یوزر.استوری} چیست؟}
  
  \باره{کاربر از این \واژه{فیچر} چه می‌خواهد؟}

  \باره{اگر سامانه این \واژه{فیچر} را داشته باشد، چه سودی به کاربر
    می‌رسد؟}

  \پایان{بارهها}

  همچنین برای هر \واژه{یوزر.استوری}، چند \واژه{اکسپتنس.کریتریا} به
  زبان مشتری تعیین می‌گردد که اگر برآورده شوند، آن \واژه{یوزر.استوری}
  به درستی محقق شده است. در \واژه{بی.دی.دی.}، \واژه{اکسپتنس.کریتریا}
  در \واژه{قالب} «سناریو» توصیف می‌شود که در بخش \رجوع{gwt} به تفصیل
  شرح داده شده است.

}

\پایان{شمارش}

\زیرقسمت{\واژه{قالب} توصیف \واژه{یوزر.استوری}‌ها و \واژه{سناریو}‌ها}
\برچسب{gwt}

سلام

\شروع{definition} به گزاره‌ای که در صورت صحّت پیاده‌سازی برقرار باشد یک
\واژه{سناریو} می‌گوییم.  \پایان{definition} \شروع{definition}

در \واژه{بی.دی.دی.} یک \واژه{سناریو} به صورت زیر نمایش
داده می‌شود:

{\نقل
\متن‌سیاه{\واژه{given}} شرط $g$ برقرار باشد، \\
\متن‌سیاه{\واژه{when}} اتفاق $w$ رخ بدهد، \\
\متن‌سیاه{\واژه{then}} شرط $t$ باید برقرار باشد.
}

و یا به طور خلاصه:

{\نقل
\متن‌سیاه{\واژه{given}} $g$ \\
\متن‌سیاه{\واژه{when}} $w$ \\
\متن‌سیاه{\واژه{then}} $t$
}

این شیوه‌ی بیان سناریو \واژه{given-when-then} نام دارد.
\پایان{definition}

%%% Local Variables:
%%% mode: latex
%%% TeX-master: "../main"
%%% End:


\قسمت{شیوه‌ی نگارش پایان‌نامه}\برچسب{sec:intro:writing}

این بخش به شیوه‌ی نگارش پایان‌نامه و چگونگی به‌کارگیری این \واژه{قالب} اختصاص دارد.
در این قالب، هر فصل از پایان‌نامه با یک \واژه{پوشه} مشخص می‌شود. برای نمونه، فصل مقدمه که در حال خواندن آن هستید، در پوشه‌ای با نام \نام{اینتروداکشِن}{introduction} نوشته شده است. در هر پوشه، یک \واژه{پرونده} با همان نام قرار دارد که متن آغازین فصل در آن نوشته می‌شود. این توصیف باید هدف از نگارش آن فصل را بیان کرده، به همه‌ی بخش‌های آن ارجاع نموده و در انتها پرونده‌های جداگانه‌ای را که به هر بخش (در همان پوشه) اختصاص پیدا کرده‌اند، با دستور \نام{اینپوت}{input} دربرگیرد.

انگیزه‌ی اصلی از تقسیم‌بندی فصل‌ها و بخش‌ها (به ترتیب) در پوشه‌ها و پرونده‌ها، این است که مدیریت حجم انبوهی از نوشته‌ها در یک پایان‌نامه، با این تقسیم‌بندی ساده‌تر خواهد شد. همچنین توصیه می‌شود از \ترجمه[سامانه‌ی پایش  نسخه‌های]{سامانه‌ی پایش  نسخه‌ها}{Version Control System} \نام{گیت}{Git} برای دنبال کردن تغییرها در متن پایان‌نامه بهره بگیرید.

اگر به متن \نام{تِک}{\TeX} این نوشته تا به اینجا نگاه کنید، چند دستور کلیدی را مشاهده می‌کنید که برای ترجمه‌ی واژه‌ها به‌کار رفته‌اند. شایان ذکر است که در همه‌ی متن پایان‌نامه، نباید از الفبای غیر پارسی (از جمله انگلیسی) بهره گرفته شود. در صورتی که می‌خواهید واژه‌ی بیگانه‌ای را به‌کار بگیرید و معادل پارسی برای آن موجود نیست، لازم است که آن واژه را با الفبای پارسی نوشته و آنگاه از زیرنویس برای بیان شیوه‌ی نگارش آن در زبان انگلیسی (برای نمونه \نام{اِسمیت}{Smith}) بهره  بگیرید. در ادامه ۴ دستور کلیدی در این رابطه معرفی می‌شوند.

\شروع{شمارش}
	\باره دستور \موکد{ترجمه} با یک آرگومان اختیاری و دو آرگومان اجباری فراخوانی می‌شود. دو آرگومان اجباری به ترتیب واژه‌ی پارسی و انگلیسی را مشخص می‌کنند. واژه‌ی پارسی در متن نوشته شده، واژه‌ی انگلیسی در زیرنویس آورده شده و دو سطر به واژه‌نامه‌ها در انتهای پایان‌نامه افزوده می‌گردد (یک سطر به واژه‌نامه‌ی پارسی به انگلیسی و یک سطر به واژه‌نامه‌ی انگلیسی به پارسی). ممکن است واژه‌ای که می‌خواهید در متن نوشته شود با واژه‌ای که می‌خواهید در واژه‌نامه آورده شود، متفاوت باشد. برای نمونه می‌خواهید واژه‌ی \ترجمه[داده‌ساختارهای]{داده‌ساختار}{Data Structure} گوناگون را در متن بیاورید، ولی حالت مفرد آن، یعنی داده‌ساختار را به واژه‌نامه اضافه کنید. در این حالت می‌توانید آنچه را که در متن نوشته می‌شود، با آرگومان اختیاری دستور \موکد{ترجمه} تعیین کنید.
	
	\باره دستور \موکد{نوواژه} همانند دستور بالا دو آرگومان اجباری برای مشخص کردن واژه‌های پارسی و انگلیسی در دو واژه‌نامه دریافت می‌کند. با این تفاوت که هیچ واژه‌ای را در متن نمی‌نویسد و به طور خودکار سطری به واژه‌نامه‌ها نمی‌افزاید. بلکه همانند تعریف مرجع‌ها در \نام{بیبتِک}{BibTeX} تنها واژه‌ها را آماده‌ی به‌کارگیری می‌کند. اگر واژه توسط دستور \موکد{واژه} در متن به‌کارگرفته شد، آن را به واژه‌نامه‌ها می‌افزاید و در غیر این صورت هیچ کاری نمی‌کند. دستورهای \موکد{نوواژه} را در پرونده‌ی خاص واژه‌نامه\زیرنویس{این پرونده در پوشه‌ی general و با نام \چر{glossaries.tex} جای دارد. به چگونگی استفاده از زیرنویس به جای به‌کارگیری الفبای انگلیسی در متن و همچنین دستور \موکد{چر} توجه کنید.} بنویسید. این کار سبب می‌شود که بتوان واژه‌نامه و جای‌گزین‌های تعریف شده توسط شما را در آینده به اشتراک گذاشت. همچنین این دستور یک آرگومان اختیاری به عنوان کلید دریافت می‌کند. در صورت عدم تعیین کلید، واژه‌ی پارسی نقش کلید را ایفا می‌کند.
	
	\باره دستور \موکد{واژه} همراه با دستور بالا کار می‌کند. این دستور کلید مشخص شده در دستور \موکد{نوواژه} را به عنوان تنها آرگومان اجباری خود دریافت کرده و واژه‌ی پارسی مرتبط را نشان می‌دهد. اگر نخستین باری باشد که آن واژه در متن آورده شده است، انگلیسی آن هم در زیرنویس آورده می‌شود. اگر همانند حالت دستور \موکد{ترجمه} بخواهید گونه‌ای متفاوت از واژه در متن آورده شود، آن را به عنوان آرگومان اختیاری دستور \موکد{واژه} مشخص کنید. توصیه می‌شود به جای \موکد{ترجمه} از \موکد{نوواژه} و \موکد{واژه} استفاده شود، زیرا در صورت جابه‌جایی بخش‌ها، تشخیص نخستین کاربرد یک واژه، خودکار انجام می‌شود، ولی در \موکد{ترجمه} باید نخستین کاربرد با دستور \موکد{ترجمه} معلوم شود و کاربردهای بعدی به صورت معمولی (بدون دستور) نوشته شوند.
	
	\باره دستور \موکد{نام} دو آرگومان اجباری را دریافت کرده و آرگومان پارسی تخست را در متن و آرگومان انگلیسی دوم را در زیرنویس  می‌آورد. این دستور به واژه‌نامه، هیچ سطری را اضافه نمی‌کند. کاربرد این دستور برای نام‌های غیر پارسی است که به دلیل مجاز نبدون به‌کارگیری الفبای انگلیسی در متن پایان‌نامه باید در زیرنویس آورده شوند.
\پایان{شمارش}


\begin{figure}[tbp]
\centering
\includegraphics[width=0.7\linewidth]{ce_department}
\def\mytempcaption{تصویری از دانشکده‌ی مهندسی کامپیوتر در دانشگاه صنعتی شریف}
\caption[\mytempcaption]{\mytempcaption~\cite{sharif:cesite}}
\label{fig:cedepartment}
\end{figure}

در ادامه نمونه‌هایی از لیست غیر شمارشی، شکل و جدول آورده شده‌اند. برای نمایش شکل‌ها، کافی است پرونده‌ی تصویر مورد نظر را در پوشه‌ی تصویرها\زیرنویس{به نام img} قرار داده و بدون اشاره به پوشه‌ی دربرگیرنده‌ی آن، تنها نام تصویر را ذکر کنید. به هر شکل و جدول باید دست کم یک بار در متن ارجاع داده شود. برای نمونه شکل~\رجوع{fig:cedepartment} تصویری از دانشکده‌ی مهندسی کامپیوتر را نشان می‌دهد. ارجاع مناسب شکل‌ها باید در عنوان آن‌ها صورت گیرد. برای اینکه این ارجاع در فهرست شکل‌ها آورده نشود از آرگومان اختیاری مرتبط همچون نمونه‌ی آورده شده، استفاده کنید.


\begin{table}[tbp]
\شرح{چهار دستور کلیدی در ترجمه‌ی واژه‌ها}
\برچسب{tbl:translate}
\centering
\begin{tabular}{ccC{7cm}}
\toprule
%\hline \rowcolor[gray]{.7}
نام دستور & نمونه & شرح \\
\midrule
ترجمه & \textbackslash ترجمه[نمونه‌های]\{نمونه\}\{Example\} & آرگومان اختیاری یا واژه‌ی پارسی را در متن آورده، واژه‌ی انگلیسی را در زیرنویس آورده و دو نگاشت را به دو واژه‌نامه خواهد افزود. \\
نوواژه & \textbackslash نوواژه[schema]\{شِما\}\{Schema\} & آرگومان اختیاری کلید را مشخص می‌کند. دو آرگومان دیگر معلوم می‌کنند در صورت به‌کارگیری کلید در متن، چه سطری به واژه‌نامه‌ها افزوده شود. کلید پیش‌فرض واژه‌ی پارسی است. \\
واژه & \textbackslash واژه[شِماهایی]\{schema\} & آرگومان اختیاری تعیین می‌کند در متن چه واژه‌ای باید نوشته شود. کلید (که به طور پیش‌فرض در متن نوشته می‌شود)، سطر مرتبط از دستور نوواژه را نشان می‌دهد. \\
نام & \textbackslash نام\{\URL\}\{URL\} & آرگومان نخست در متن (نشان‌دهنده‌ی یک نام) و آرگومان دوم در زیرنویس آورده می‌شود. \\
\bottomrule
\end{tabular}
\end{table}

جدول~\رجوع{tbl:translate} برای نمونه، فهرستی از چهار دستور توصیف شده‌ی بالا را گردآوری کرده است.
برای جدول‌ها، همانند این نمونه از \واژه[بسته‌ی]{بسته} \booktabs{} بهره گرفته و از خط‌های عمودی یا رنگ کردن خانه‌ها در جدول پرهیز کنید.
برای آخرین نمونه یک لیست بدون ترتیب آورده شده است که در آن برخی از خطاهای رایج در نگارش متن‌های پارسی ذکر شده‌اند.

\شروع{بارهها}
	\باره حرف اضافه‌ی \موکد{را} هرگز نباید بعد از فعل آورده شود،
	
	\باره ترجیح بر نگارش واژ \موکد{ی} بزرگ به جای گونه‌ی کوچک آن است. برای نمونه خانه‌ی ما بر خانهٔ ما ترجیح دارد. دلیل این ترجیح اشتباه شدن گونه‌ی کوچک \موکد{ی} با واژ همزه است،
	
	\باره جمع مکسر، همزه و تنوین در پارسی وجود نداشته و باید با شکل‌های متناسب جای‌گزین شوند. می‌توانید به جای مثلاً بگویید برای نمونه، به جای مسئله بنویسید مساله (برای عدم نیاز به نگارش همزه)، به جای تغییرات بگویید تغییرها و به جای موارد بگویید موردها،
	
	\باره برای اینکه تشخیص دهید کدام واژه‌ها را باید جدا نوشت و کدام واژه‌ها را باید یکپارچه نوشت، از یک قاعده‌ی ساده پیروی نمایید:
	اگر معنی واژه‌ای از اجزای آن بدون کم و کاستی قابل برداشت است، آن را جدا (با نیم‌فاصله\زیرنویس{به نیم‌فاصله، فاصله‌ی مجازی هم می‌گویند.}) بنویسید و اگر باید آن واژه را به عنوان واژه‌ای مستقل به خاطر سپرد،
	یکپارچه (سرهم) بنویسید. برای نمونه \موکد{جوانمرد} باید یکپارچه نوشته شود، ولی \موکد{جای‌گزین} باید جدا نوشته شود،
	
	\باره اگر از \ترجمه{ سامانه‌ی عامل}{Operating System} \نام{ویندوز}{Windows} استفاده می‌کنید، صفحه‌کلید \ترجمه[استانده‌ی]{استانده}{Standard} ایران را از~\مرجع{persiancomputing:keyboard} دریافت کنید. اگر از  سامانه‌ی عامل \نام{لینوکس}{Linux} استفاده می‌کنید،   صفحه‌کلید استانده‌ی ایران  به طور پیش‌فرض پشتیبانی می‌شود.
	
\پایان{بارهها}

در پایان هر بخش یا فصل مناسب است در جمله‌ای کوتاه به بخش یا فصل بعدی اشاره کنید. در بخش بعدی ساختار این پایان‌نامه بیان می‌گردد.


\قسمت{ساختار پایان‌نامه}\برچسب{sec:intro:structure}

% write this at the end. -- 0.25 col
این بخش ساختار پایان‌نامه را مشخص می‌کند. همان‌گونه که هر پایان‌نامه‌ای با فصل مقدمه شروع شده و به فصلی درباره‌ی نتیجه‌گیری و سوی کارهای آتی آن پایان‌نامه ختم می‌شود، لازم است بخشی در پایان فصل مقدمه به توصیف ساختار کلی پایان‌نامه اختصاص پیدا کند.
این فصل باید به همه‌ی فصل‌ها ارجاع نماید و به طور خلاصه آنچه را در هر فصل بیان خواهد شد ذکر نماید. برای نمونه به بند بعدی توجه کنید.

در ادامه و پس از پایان فصل مقدمه، فصل~\رجوع{chap:future} به جمع‌بندی آنچه در  این  پایان‌نامه مورد بحث قرار گرفت پرداخته و دست‌آوردهای آن و سوی کارهای آتی آن را بیان  می‌کند.


%%% Local Variables:
%%% mode: latex
%%% TeX-master: "../main"
%%% End:
